% !TeX root = bericht.tex
\section{Aufgabe 4 - PI-Regler mit Lastmoment}
In den vorherigen Aufgaben haben wir unser Motor ohne Lastmoment geregeelt.
Hierfür müssen wir unsere Regelstrecke leicht anpassen um die Störung zu berücksichtigen.
\begin{figure}[H]
    \centering
    \includegraphics[width=1.0\textwidth]{strecke_mit_stoerung.png}
    \caption{Simulink Modell der Regelstrecke mit Stoerung}
\end{figure}

\begin{figure}[H]
    \centering
    \includegraphics[width=1.0\textwidth]{pi-regler_mit_stoerung.png}
    \caption{Simulink Modell des Reglers mit PI-Regler mit Lastmoment}
\end{figure}
In der Simulation wurde ein Lastmoment von $0.01 Nm$ bei $t=10s$ hinzugefügt.

\begin{figure}[H]
    \centering
    \includegraphics[width=0.9\textwidth]{a4_pi_regelung_sim_100rad_stoerung.png}
    \caption{Simulationsergebnis der Regelung mit einem PI-Regler auf 100 rad/s mit Lastmoment}
\end{figure}

Hier kann man gut erkennen, dass der PI-Regler auch eine Störung von $10 mNm$ ausgleichen kann.
